\documentclass[a4paper,11pt,twocolumn]{article}
\usepackage[margin=2.54cm,columnsep=0.6cm]{geometry}
\usepackage{newtxtext}
\usepackage{amsmath,amssymb}
\usepackage{booktabs}
\usepackage{graphicx}
\usepackage{natbib}
\usepackage{titlesec}
\usepackage{url}
\setlength{\parindent}{0.38cm}
\setlength{\parskip}{0pt}
\titleformat{\section}{\bfseries\fontsize{12}{14}\selectfont}{\thesection}{0.5em}{}
\titleformat{\subsection}{\bfseries\fontsize{11}{13.6}\selectfont}{\thesubsection}{0.5em}{}
\titlespacing*{\section}{0pt}{7pt}{3pt}
\titlespacing*{\subsection}{0pt}{5pt}{2pt}
\pagestyle{empty}
\title{\fontsize{14.5}{17}\selectfont\textbf{QWeave: Correctness-First Deterministic and Learned Qubit Routing}}
\author{%
\fontsize{12}{14}\selectfont
\textbf{Atharva Sheersh Pandey}$^{1}$, \textbf{Diptesh Das}$^{1}$,
\textbf{Shrivardhini N}$^{1}$,\\
\textbf{Haridasu Sreedhar}$^{1}$, \textbf{Prof. Bhuvaneswari M}$^{2}$\\[2pt]
\texttt{atharva.sheersh2024@vitstudent.ac.in};
\texttt{diptesh.das2024@vitstudent.ac.in};\\
\texttt{shrivardhini.n2024@vitstudent.ac.in};
\texttt{haridasu.sreedhar2024@vitstudent.ac.in}\\[2pt]
$^{1}$School of Computer Science and Engineering, VIT Vellore\\
$^{2}$Faculty Guide, School of Computer Science and Engineering, VIT Vellore
}
\date{}

\begin{document}
\maketitle
\begin{abstract}
Restricted two-qubit connectivity forces a compiler to place logical qubits
on physical sites and insert movement before nonlocal interactions. This paper
studies QWeave as a correctness-first compiler rather than treating low depth
alone as sufficient evidence. The deterministic path combines weighted initial
mapping, strict local improvement, and shortest-path SWAP routing; identity
mapping, Qiskit SABRE, and a small exact initial-mapping oracle are distinct
references. We prove termination of the finite mapping descent, preservation of
the routing layout invariant and supported unitary semantics, feasibility and
fixed-DAG optimality of the list scheduler, equivalence of directed-CNOT
lowering, and termination of the guarded learned route. These guarantees do not
prove policy quality. That question is evaluated on frozen benchmark v2:
180 synthetic cases, five final training seeds, and 792 held-out outputs.
Every tested output passed its applicable numerical semantic check and every
direction-lowered cost circuit passed ordered-coupler legality. The paired
GNN--PPO minus weighted depth difference was 0 gates with a 95\% bootstrap
interval of $[0,0]$; GNN--PPO and no-message-passing PPO also had paired median
depth and SWAP differences of 0. Training improved markedly on an untrained
policy, but the study supplies no evidence of a message-passing or
learned-routing advantage over the deterministic references.
\end{abstract}

\section{Introduction}
Quantum algorithms are commonly written with unrestricted two-qubit
interactions, whereas a device exposes a finite coupling graph. Compilation
must therefore choose an initial placement and transform every nonlocal
interaction into operations supported by that graph. Qiskit describes layout,
routing, translation, optimization, and scheduling as separate stages
\citep{ibmtranspiler}. SABRE uses bidirectional heuristic search for layout and
SWAP selection \citep{li2019sabre}; it is consequently a substantive baseline.

QWeave asks a narrower question: under a declared circuit language, topology,
mapping convention, and validation protocol, what do deterministic placement,
guarded learned routing, and post-routing scheduling each contribute? The
answer requires more than an edge-legality check. An emitted circuit can use
only valid couplers and still implement the wrong logical computation if a
layout update, source SWAP, parameter, or classical destination is mishandled.
For that reason, QWeave combines trace replay with numerical checks and keeps
mapping, routing, scheduling, and direction lowering as separately reported
stages.

This paper contributes (i) an explicit deterministic compiler and adapter
contract, (ii) formal correctness arguments for the invariants that can be
proved from that contract, (iii) a guarded GNN--PPO implementation with
message-passing and training controls, and (iv) a frozen held-out comparison
with paired uncertainty estimates and counterexamples. Learning-based routing
already exists \citep{pozzi2020rl,sinha2022gnn}; neither the use of PPO nor a
graph encoder is claimed as novel. The result is also deliberately bounded:
the tested policy does not outperform the weighted deterministic reference.

\section{Related Work}
Earlier qubit-allocation systems model placement and movement under
architecture constraints. Siraichi et al. formulate qubit allocation as a
compiler problem, while Zulehner et al. combine architecture-aware mapping with
heuristic search for IBM devices \citep{siraichi2018allocation,zulehner2019efficient}.
Cowtan et al. give an architecture-agnostic routing methodology, and Childs et
al. study circuit transformations through routing via matchings and token
swapping \citep{cowtan2019routing,childs2019transformations}. These works
motivate separating initial placement from dynamic movement.

SABRE combines reverse traversal for initial layout with heuristic SWAP search
\citep{li2019sabre}. In Qiskit, its layout pass may also invoke routing and add
idle physical qubits \citep{ibmsabre}. Pozzi et al. optimize routing with deep
reinforcement learning, while Sinha et al. combine a graph neural network with
Monte Carlo tree search \citep{pozzi2020rl,sinha2022gnn}. QWeave instead uses
a masked graph actor--critic trained by PPO, drawing on graph convolution,
clipped policy optimization, and generalized advantage estimation
\citep{kipf2017gcn,schulman2017ppo,schulman2015gae}.

Exact optimization provides a different reference. Nannicini et al. formulate
joint assignment and routing by integer programming
\citep{nannicini2022ip}. QWeave's CP-SAT component is intentionally smaller:
it optimizes only a static initial-mapping surrogate for at most eight logical
qubits. Its optimal status cannot certify minimum SWAP count or compiled depth.

\section{Problem Formulation and Deterministic Path}
Let $C=(g_1,\ldots,g_T)$ be an ordered supported circuit on logical qubits
$L=\{0,\ldots,n-1\}$ and let $H=(P,E_H)$ be a connected undirected routing
graph. An initial mapping $m_0:L\rightarrow P$ is injective and uses the
convention $m(\mathrm{logical})=\mathrm{physical}$. The interaction graph
$G_C$ includes isolated logical qubits. An edge $(i,j)$ has weight $w_{ij}$
equal to its two-qubit interaction count. QWeave minimizes the placement
surrogate
\begin{equation}
 J(m)=\sum_{(i,j)\in E_C}w_{ij}\,d_H(m(i),m(j)),
 \label{eq:objective}
\end{equation}
where $d_H$ is shortest-path distance. Optional temporal decay replaces an
interaction count by $\delta^t$, but benchmark v2 uses integer counts.

The mapper places a high weighted-degree logical qubit on a high-centrality
physical node, then greedily assigns the remaining logical qubits using
interaction weight and partial distance cost. Deterministic ordering resolves
ties. Pairwise local search evaluates exchanges of occupied physical sites and
accepts the best exchange only when it strictly decreases $J$.

For each source operation, the deterministic router emits an adjacent
one- or two-qubit gate immediately. If a two-qubit gate is nonlocal, it selects
a deterministic shortest path and moves the first logical qubit toward the
second through edge-local SWAPs. Both forward and inverse layouts are updated
after every inserted SWAP. The Basic reference supplies identity mapping to
this router; weighted mapping supplies the greedy local-search layout. SABRE is
a separate Qiskit compiler path with a fixed seed.

The Qiskit adapter accepts unconditioned one- and two-qubit gates, one-qubit
resets, and a terminal measurement suffix while preserving global phase and
classical destinations. It rejects barriers, higher-arity gates, dynamic
control flow, and classically conditioned operations instead of silently
passing unsupported semantics. A trace records one event for each source gate
or inserted SWAP. Replay checks source order, parameters, targets, layout
updates, and the final mapping. SABRE receives both orientations of every
undirected routing edge so all methods share the same adjacency interpretation.

\section{Formal Correctness Arguments}
The propositions below cover algorithmic invariants under the stated input
contract. They do not establish that the static objective predicts minimum
route cost, that a learned policy is good, or that finite numerical probes are
a proof over all inputs.

\paragraph{Proposition 1 (mapping-search termination).}
The strict pair-exchange local search terminates after finitely many accepted
moves.

\paragraph{Proof.}
There are $|P|!/(|P|-|L|)!$ injective maps from $L$ to $P$, so the search space
is finite. Every accepted move produces a previously unvisited objective level
with strictly smaller $J$. A strictly descending sequence over a finite set
cannot be infinite. The implementation's iteration cap is therefore an
additional guard, not the basis of termination.

\paragraph{Proposition 2 (layout and legality invariant).}
After every emitted operation, the forward layout remains injective, its
inverse agrees on occupied sites, and every emitted two-qubit operation uses an
edge of $E_H$.

\paragraph{Proof.}
The claim holds initially by the mapping contract. An inserted SWAP exchanges
the logical occupants of two adjacent physical sites; updating both maps by
the same transposition preserves injectivity and inverse agreement. A source
one-qubit gate cannot violate edge legality. A source two-qubit gate is emitted
only when its current images are adjacent, either immediately or after the
shortest-path swaps. Induction over the emitted sequence proves the claim.

\paragraph{Theorem 1 (supported unitary semantics).}
For an equal-width circuit containing only supported unitary one- and
two-qubit gates, let $R_k$ be the physical route after the first $k$
source gates, let $C_k$ be the corresponding logical prefix, and let $m_k$
be the current layout. With $P_m$ denoting the wire
permutation induced by $m$,
\begin{equation}
 U_{R_k}P_{m_0}=P_{m_k}U_{C_k}
 \quad\text{for }0\leq k\leq T .
 \label{eq:route-invariant}
\end{equation}

\paragraph{Proof.}
At $k=0$ both prefixes are identities, so the equality holds. Assume it holds
through source gate $k$. Every inserted physical SWAP permutes two adjacent
wires, and the layout update records exactly that same transposition. The
physical state interpreted through the updated layout therefore represents the
same logical state. Once the operands are adjacent, applying the source gate to
their current physical images has the same action as applying the logical gate
before the current wire permutation. Composing these steps with the induction
hypothesis proves the invariant for $k+1$. At $k=T$, the final routed unitary
equals the logical unitary up to the declared initial and final permutations.

\paragraph{Proposition 3 (fixed-route scheduling).}
For positive operation durations and no resource constraints beyond physical
qubit exclusivity, the list scheduler returns a feasible schedule. Its
makespan equals the weighted critical-path length of the fixed per-qubit
precedence DAG and is therefore minimal among schedules that preserve that DAG.

\paragraph{Proof.}
The DAG contains an edge between successive operations on each physical qubit.
Consequently, incomparable ready operations share no qubit and may overlap.
For operation $i$, the scheduler uses
\begin{equation}
 r_i=d_i+\max_{j\in\mathrm{succ}(i)}r_j,
 \qquad
 s_i=\max_{j\in\mathrm{pred}(i)}(s_j+d_j),
 \label{eq:schedule}
\end{equation}
with an empty maximum equal to zero. Thus each operation starts after every
predecessor finishes, proving feasibility. Every schedule respecting the DAG
has makespan at least its weighted critical path. Earliest-start recursion
attains that bound, so the returned makespan is optimal for this fixed DAG.
This statement does not allow gate commutation or model control electronics.

\paragraph{Proposition 4 (directed lowering).}
If at least one orientation of an undirected routing edge supports CNOT, the
lowering pass implements either requested CNOT orientation and an edge-local
SWAP using only declared directed CNOTs and one-qubit Hadamards.

\paragraph{Proof.}
For a supported physical CNOT directed from $v$ to $u$, the opposite
orientation satisfies
\begin{equation}
 \mathrm{CX}_{u\rightarrow v}
 =(H_u\otimes H_v)\,
 \mathrm{CX}_{v\rightarrow u}\,
 (H_u\otimes H_v).
 \label{eq:reverse-cx}
\end{equation}
An undirected SWAP is three alternating CNOTs. Replacing each unavailable
orientation by Equation~\eqref{eq:reverse-cx} preserves the product unitary
and leaves every two-qubit instruction on an allowed ordered arc.

\paragraph{Proposition 5 (guarded-route termination).}
For a finite supported circuit on a finite connected routing graph, the
guarded learned environment terminates and returns a legal route.

\paragraph{Proof.}
Each policy-controlled segment has a finite step limit. On repeated
nonprogress or exhaustion, fallback selects a finite shortest path, performs
adjacent swaps, and executes the current gate. Hence fallback advances the
source index. A finite circuit permits only finitely many such advances, and
Proposition 2 gives legality for every fallback operation. Replay validation
must succeed before the environment returns a result.

\section{Learned Router}
The Gymnasium environment exposes an \textsc{Execute} action and one SWAP
action for every sorted hardware edge. Its observation contains five features
per physical node, four per hardware edge, a five-value frontier description,
the current logical-to-physical mapping, and a binary legal-action mask.
Illegal actions do not mutate the circuit or layout. Fallback counts are
reported because the termination guard can conceal a weak policy.

The actor--critic performs two rounds of edge-conditioned message passing and
scores every edge SWAP plus \textsc{Execute}. PPO uses clipped policy ratios,
generalized advantage estimation, value and entropy terms, and gradient
clipping \citep{schulman2017ppo,schulman2015gae}. A zero-message-passing PPO
control preserves the optimizer and reward, while an untrained graph control
preserves the architecture without PPO updates. All variants start from the
weighted deterministic map, making this a routing-policy comparison rather
than a learned-layout experiment.

\section{Validation and Experimental Protocol}
Structural replay checks mapping completeness, injectivity, source-gate order,
inserted SWAP consistency, global phase, and edge legality. For equal-width
unitary circuits of at most six qubits, QWeave evaluates
Equation~\eqref{eq:route-invariant} at $k=T$ as an exact operator equality.
Wider unitary circuits and outputs with idle physical sites use reproducible
layout-aware statevector probes up to 12 physical qubits. Exact equality is a
semantic proof for the evaluated circuit; a finite probe set is evidence that
can detect faults but is not a universal proof. Reset and terminal-measurement
circuits receive structural replay. Branch-dependent dynamic control remains
unsupported.

After routing, the fixed-route scheduler measures depth or declared-duration
makespan on the same legal route. A separate hardware-cost pass lowers CNOTs
and SWAPs to declared directed couplers, validates ordered arcs, then evaluates
the schedule and an independent-error proxy. Durations and edge errors are
synthetic assumptions rather than device calibrations.

Benchmark version 2.0 freezes 180 cases at four, six, and eight logical
qubits, with two instances per family. Sixty training cases use GHZ,
layered-CX, seeded-interaction, Bernstein--Vazirani, and QAOA MaxCut families;
the latter follows the alternating-operator construction of
\citet{farhi2014qaoa}. Eighty-four validation cases test topology transfer and
held-out families. The 36 untouched test cases combine QFT circuits, motivated
by the approximate transform of \citet{coppersmith2002qft}, and
hardware-efficient circuits with two-row grids, chorded rings, and larger
grids containing idle sites. Every case stores QASM, source hash, undirected
and directed edges, and declared synthetic error values.

Three candidates receive 1,024 training steps and validation seed 101. The
selection score is median depth plus half the SWAP count plus five times the
fallback count. The selected candidate is trained for 4,096 steps at seeds 7,
17, 29, 41, and 53 for both GNN--PPO and the no-message control. The untrained
control uses the same initialization seeds. SABRE uses the same five seeds;
Basic and weighted routing are deterministic. Runtime is the median of three
measurements after one warmup. Seeds are collapsed within each case before
paired differences are formed. Confidence intervals use 5,000 deterministic
bootstrap resamples with seed 20260921.

\section{Results}
A fresh local run passed 69 tests under Python 3.11.15 and Qiskit 2.5.2. The
aggressive validation candidate scored 19.5; conservative and standard each
scored 24.0. Table~\ref{tab:test-results} summarizes the 36 held-out cases.
For seeded methods, each case contributes its within-case seed median. All 792
outputs passed their applicable numerical check. QWeave routes also passed
symbolic replay; SABRE has no QWeave trace. Every unit-duration scheduling
delta was zero, as expected when circuit depth already expresses the same
per-qubit precedence.

\begin{table*}[t]
\centering
\caption{Benchmark-v2 test summary. Runtime is median wall time in seconds.}
\label{tab:test-results}
\small
\begin{tabular}{lrrrr}
\toprule
Method & Depth & SWAP & Runtime & Fallback \\
\midrule
Basic & 48.5 & 11.5 & 0.0022 & 0 \\
Weighted & 36.0 & 0.0 & 0.0036 & 0 \\
SABRE & 46.0 & 7.0 & 0.0028 & 0 \\
GNN--PPO & 37.5 & 0.0 & 0.3473 & 4543 \\
No-message PPO & 36.0 & 0.0 & 0.1040 & 742 \\
Untrained GNN & 198.0 & 193.0 & 0.6681 & 8059 \\
\bottomrule
\end{tabular}
\end{table*}

The paired GNN--PPO minus weighted depth difference was 0 gates with interval
$[0,0]$; versus SABRE it was 0 with interval $[0,5]$. Its paired SWAP
difference versus weighted routing was 0 with interval $[0,5]$. GNN--PPO and
no-message PPO had paired median depth and SWAP differences of 0 with
intervals $[0,0]$, so the experiment supplies no evidence that message passing
helped. Training did matter relative to the untrained GNN: the paired depth
difference was $-162.5$ with interval $[-184,-106]$, and the SWAP difference
was $-175$ with interval $[-223,-148]$.

Family results explain the aggregate. On hardware-efficient circuits,
weighted, GNN--PPO, and no-message PPO all had median depth 27 and no SWAPs. On
QFT, weighted routing had median depth 71.5 and 16.5 SWAPs, GNN--PPO had 83.5
and 41, and no-message PPO had 90 and 31.5. GNN--PPO improved by eight depth
layers on \texttt{qft\_i1\_6q\_chorded\_ring}, but regressed by 89 depth gates
and 130 SWAPs on \texttt{qft\_i1\_8q\_grid2}. It invoked 4,543 test
fallbacks, compared with 742 for no-message PPO and zero for deterministic
methods. Seed 29 was unstable, ending with 212 training fallbacks and a recent
mean return near $-103$.

\begin{figure*}[t]
\centering
\includegraphics[width=0.94\textwidth]{../results/benchmark_v2/figures/learning_curves.png}
\caption{Five-seed learning curves. Lines show the median recent episode
return and bands show the interquartile range. Neither curve is monotone, and
the GNN curve degrades more strongly late in training.}
\label{fig:learning-curves}
\end{figure*}

Before direction lowering, most routes were illegal on the one-way couplers.
After reverse-CNOT and SWAP synthesis, all 792 cost circuits passed ordered-arc
legality. Median added depth was 32 for weighted routing, 34 for SABRE, 37.5
for GNN--PPO, 28.5 for no-message PPO, 57 for Basic, and 441 for the untrained
GNN. The corrected paired GNN--PPO minus weighted makespan difference was
0 ns with interval $[0,1113]$. These synthetic proxy costs provide no learned
advantage.

\section{Threats to Validity}
The formal results depend on the supported gate contract, a connected routing
graph, correct primitive gate implementations, and the absence of unmodelled
shared hardware resources. They prove preservation properties of the
algorithms, not software free of defects; tests and independent reruns address
implementation risk. Numerical probes for wider circuits are incomplete by
construction. Measurement, reset, and dynamic-control semantics do not receive
the same exact unitary argument.

The empirical study uses synthetic circuits, durations, and error rates. Its
five seeds and 36 held-out cases support paired statements about benchmark v2,
not general performance or state-of-the-art rank. Runtime is machine-specific,
and fallback makes the learned methods safe while also weakening the
interpretation of policy-only behavior. The retained pre-correction record is
disclosed because its duration fields scheduled pre-lowering circuits; the
schema-2.1 replay recomputed only direction-dependent cost fields from saved
checkpoints and frozen cases.

\section{Conclusion}
QWeave combines an explicit deterministic reference, a small mapping oracle,
fixed-route scheduling, directed lowering, and a guarded learned router.
Formal arguments now state why strict mapping descent terminates, why routing
preserves the declared layout semantics, why the fixed-DAG scheduler is
feasible and makespan-optimal under its model, why direction reversal is
unitarily sound, and why fallback terminates. Benchmark v2 complements those
proofs with held-out evidence. PPO training clearly improves on an untrained
policy, but the tested message-passing component and learned route do not
improve on weighted routing. Larger preregistered device-derived studies are
needed before any superiority or calibrated-fidelity claim.

\bibliographystyle{plainnat}
\bibliography{references}
\end{document}

